\documentclass[12pt,a4paper]{article}
\usepackage[T1]{fontenc}
\usepackage[scaled]{uarial}
\renewcommand*\familydefault{\sfdefault}
\usepackage{longtable}

\begin{document}
	\section{Daitch-Mokotoff Soundex Coding}
	
	The Daitch-Mokotoff Soundex System was created by Randy Daitch and Gary 
	Mokotoff of the Jewish Genealogical Society (New York), because they 
	concluded the system developed by Robert Russell in 1918 in use today by 
	the U.S. National Archives and Records Administration (NARA) does not apply
	well to many Slavic and Yiddish surnames.  Daitch-Mokotoff Soundex also 
	includes refinements that are independent of ethnic considerations.  The 
	rules for converting surnames into D-M Code numbers are listed below.  They
	are followed by the coding chart. 
	
	\begin{enumerate}
	\item Names are coded to six digits, each digit representing a sound listed 
	in the coding chart (below).
	
	\item When a name lacks enough coded sounds for six digits, use zeros to 
	fill to six digits.  GOLDEN which has only four coded sounds [G-L-D-N]
	is coded as 583600.
	
	\item The letters A, E, I, O, U, J, and Y are always coded at the beginning 
	of a name as in Alpert 087930.  In any other situation, they are 
	ignored except when two of them form a pair and the pair comes before 
	a vowel, as in Breuer 791900 but not Freud.
	
	\item The letter H is coded at the beginning of a name, as in Haber 579000, 
	or preceding a vowel, as in Manheim 665600, otherwise it is not coded.
	
	\item When adjacent sounds can combine to form a larger sound, they are 
	given the code number of the larger sound.  Mintz which is not coded 
	MIN-T-Z but MIN-TZ 664000. 
	
	\item When adjacent letters have the same code number, they are coded as one
	sound, as in TOPF, which is not coded TO-P-F 377000 but TO-PF 370000. 
	Exceptions to this rule are the letter combinations MN and NM, whose 
	letters are coded separately, as in Kleinman, which is coded 586660 
	not 586600. 
	\item When a surname consists or more than one word, it is coded as if one 
	word, such as "Ben Aron", which is treated as "Benaron". 
	
	\item Several letter and letter combinations pose the problem that they may 
	sound in one of two ways.  The letter and letter combinations CH, CK, 
	C, J, and RS are assigned two possible code numbers. 
	\end{enumerate}
	
	\section{The Daitch-Mokotoff Soundex Coding Chart}
	\begin{longtable}{|l|l|l|l|l|}
		\hline
		Letter & Alternatives & At Start & Before vowel & Any other pos \\
		\hline
		~&~&~&~&~\\
		AI & AJ, AY & 0 & 1 & NC \\
		AU &   & 0 & 7 & NC \\
		Ą & (Polish a-ogonek) & NC & NC & 6 or NC \\
		A &   & 0 & NC & NC \\
		B &   & 7 & 7 & 7 \\
		CHS &   & 5 & 54 & 54 \\
		CZ & CS, CSZ, CZS & 4 & 4 & 4 \\
		DRZ & DRS & 4 & 4 & 4 \\
		DS & DSH, DSZ & 4 & 4 & 4 \\
		DZ & DZH, DZS & 4 & 4 & 4 \\
		D & DT & 3 & 3 & 3 \\
		EI & EJ, EY & 0 & 1 & NC \\
		EU &   & 1 & 1 & NC \\
		Ę & (Polish e-ogonek) & NC & NC & 6 or NC \\
		E &   & 0 & NC & NC \\
		FB &   & 7 & 7 & 7 \\
		F &   & 7 & 7 & 7 \\
		G &   & 5 & 5 & 5 \\
		H &   & 5 & 5 & NC \\
		IA & IE, IO, IU & 1 & NC & NC \\
		I &   & 0 & NC & NC \\
		KS &   & 5 & 54 & 54 \\
		KH &   & 5 & 5 & 5 \\
		K &   & 5 & 5 & 5 \\
		L &   & 8 & 8 & 8 \\
		MN &   &   & 66 & 66 \\
		M &   & 6 & 6 & 6 \\
		NM &   &   & 66 & 66 \\
		N &   & 6 & 6 & 6 \\
		OI & OJ, OY & 0 & 1 & NC \\
		O &   & 0 & NC & NC \\
		P & PF, PH & 7 & 7 & 7 \\
		Q &   & 5 & 5 & 5 \\
		R &   & 9 & 9 & 9 \\
		SCHTSCH & SCHTSH, SCHTCH & 2 & 4 & 4 \\
		SCH &   & 4 & 4 & 4 \\
		SHTCH & SHCH, SHTSH & 2 & 4 & 4 \\
		SHT & SCHT, SCHD & 2 & 43 & 43 \\
		SH &   & 4 & 4 & 4 \\
		STCH & STSCH, SC & 2 & 4 & 4 \\
		STRZ & STRS, STSH & 2 & 4 & 4 \\
		ST &   & 2 & 43 & 43 \\
		SZCZ & SZCS & 2 & 4 & 4 \\
		SZT & SHD, SZD, SD & 2 & 43 & 43 \\
		SZ &   & 4 & 4 & 4 \\
		S &   & 4 & 4 & 4 \\
		TCH & TTCH, TTSCH & 4 & 4 & 4 \\
		TH &   & 3 & 3 & 3 \\
		TRZ & TRS & 4 & 4 & 4 \\
		TSCH & TSH & 4 & 4 & 4 \\
		TS & TTS, TTSZ, TC & 4 & 4 & 4 \\
		TZ & TTZ, TZS, TSZ & 4 & 4 & 4 \\
		Ţ & (Romanian t-cedilla) & 3 or 4 & 3 or 4 & 3 or 4 \\
		T &   & 3 & 3 & 3 \\
		UI & UJ, UY & 0 & 1 & NC \\
		U & UE & 0 & NC & NC \\
		V &   & 7 & 7 & 7 \\
		W &   & 7 & 7 & 7 \\
		X &   & 5 & 54 & 54 \\
		Y &   & 1 & NC & NC \\
		ZDZ & ZDZH, ZHDZH & 2 & 4 & 4 \\
		ZD & ZHD & 2 & 43 & 43 \\
		ZH & ZS, ZSCH, ZSH & 4 & 4 & 4 \\
		Z &   & 4 & 4 & 4 \\
		\hline
	\end{longtable}
	
		 \newpage
		 
	\section{Examples of Daitch-Mokotoff Soundex Coding}
	\paragraph{AUERBACH = 097500\\} 
	~\\
	\begin{tabular}{|l|l|l|l|l|l|l|}
		\hline
		A & UE & R & B & A & CH ~ &\\
		\hline
		0 & NC & 9 & 7 & NC & 5 & Pad \\
		\hline
		0  & ~ & 9 & 7 & &   5 & 00	\\
		\hline
	\end{tabular}

	 
	\paragraph{OHRBACH = 097500\\} 
	~\\
	\begin{tabular}{|l|l|l|l|l|l|l|}
		\hline
		O & H & R & B & A & CH ~ &\\
		\hline
		0 & NC & 9 & 7 & NC & 5 & Pad \\
		\hline
		0  & ~ & 9 & 7 & &   5 & 00	\\
		\hline
	\end{tabular}
	
	\paragraph{LIPSHITZ = 874400\\} 
	~\\
	\begin{tabular}{|l|l|l|l|l|l|l|}
		\hline
		L & I & P & SH & I & TZ ~ &\\
		\hline
		8 & NC & 7 & 4 & NC & 4 & Pad \\
		\hline
		8  & ~ &7  & 4 & &   4 & 00	\\
		\hline
	\end{tabular}
	
	\paragraph{LIPPSZYC = 874400\\} 
	~\\
	\begin{tabular}{|l|l|l|l|l|l|l|l|}
		\hline
		L & I & P & P & SZ & Y & C & ~ \\
		\hline
		8 & NC & 7 & NC & 4 & NC & 4 & Pad \\
		\hline
		8 &   & 7 &   & 4 & NC & 4 & 00 \\
		\hline
	\end{tabular}
	
	\paragraph{LEWINSKY = 876450\\}
	
	\paragraph{LEVINSKI = 876450\\}
	
	\paragraph{SZLAMAWICZ = 486740\\}
	
	\paragraph{SHLAMOVITZ = 486740\\}
	
\end{document}